\section{Conclusion}\label{sec:conclusion}
% =====================================================================
% BAN NHAP (draft) — RW (Lộc). Cau tra loi RQ1-RQ3 dang la KHUNG cho:
% dien so that sau khi do xong phia LLM (cho Hai sinh test).
% =====================================================================

We evaluated GPT-4o-mini as a unit-test generator for 120 real-world
Java and Python functions of medium cyclomatic complexity (CC 5--10),
measured inside the original projects and compared against EvoSuite
and Randoop under equal 60-second budgets.

\textbf{RQ1 (coverage threshold).} GPT-4o-mini does \emph{not} reach a median branch
coverage $\geq 80\%$, at any level of analysis. Counting all 120 sampled functions the
median is 0\% (one-sample Wilcoxon, $H0$ not rejected); even restricted to the 28 functions
whose suites compiled \emph{and} exercised the target, the median is 75.0\% ($p=0.84$,
$r=-0.22$) --- still short of 80\%. For context, EvoSuite itself only reaches a median 55.0\%
branch coverage on the same 60 Java functions and Randoop 21.1\%, so the 80\% bar is
demanding for every tool on this CC~5--10 sample, not uniquely hard for the LLM.

\textbf{RQ2 (quality vs.\ baselines).} Mutation scores are weaker still: the effective-subset
median is 36.84\% for Python and 0\% for Java, both far below the 60\% threshold. Paired
Wilcoxon tests against the search-based Java baselines show GPT-4o-mini losing decisively ---
vs.\ EvoSuite $p<0.001$, $r=-1.00$; vs.\ Randoop $p=0.003$, $r=-1.00$, the maximum possible
effect size, meaning the baseline wins every function where the two differ. Against the
Python baseline Pynguin the comparison reverses ($p=0.010$, $r=+0.85$ in GPT-4o-mini's
favour), but only because Pynguin itself achieved 0.0\% median coverage on this dataset ---
a tooling failure (Section~\ref{sec:threats}), not evidence that GPT-4o-mini's tests are
strong in absolute terms.

\textbf{RQ3 (complexity correlation).} We find no significant correlation between cyclomatic
complexity and coverage within the CC~5--10 band we sampled, in either language (Python
$\rho=+0.18$, $p=0.41$; Java $\rho=+0.13$, $p=0.83$; pooled $\rho=+0.16$, $p=0.40$). This does
\emph{not} mean complexity is irrelevant: it means that within this narrow, pre-selected
medium-complexity band, quality is dominated by the binary compiles/touches-target outcome
(Section~\ref{sec:results}) rather than by a graded decline with CC. Locating the true
``breaking point'' would require sampling a wider CC range, including low- and
high-complexity functions for contrast.

Beyond the numeric answers, two lessons generalise. First, evaluating
LLM-generated tests on \emph{extracted} functions systematically
understates the difficulty of real code: functions mined from mature
projects depend on their module context, and measurement is only
meaningful inside that context. Second, the measurement pipeline
itself is a first-class threat: most of the failures we encountered
were silent (exit code 0), and without artefact-level verification
they would have produced confident but wrong numbers.

\textbf{Future work.} Two exploratory, post-hoc follow-ups (RQ4 and RQ5,
Sections~\ref{sec:rq4}--\ref{sec:rq5}) tested whether low-cost prompting interventions close
the gap to the pre-registered thresholds. They do not, fully: Java's effective-suite rate
rises across three stages (8.3\% $\to$ 18.3\% $\to$ 25.0\% with, respectively, no
intervention, real API context, and one feedback/repair iteration on top of context) with a
maximal effect size but without reaching $p<0.05$ at $n=60$; mutation score stays flat at
0\% throughout. Python shows no benefit from context and an actual regression from
single-shot repair, most plausibly because visible failures prompt gpt-4o-mini to
\emph{expand} the suite rather than narrowly fix the failing assertion. Per our own
anti-HARKing commitment, we ran exactly these two follow-ups, reported both outcomes as
found, and stopped rather than searching further for a technique that crosses the threshold.
Useful directions for future, properly pre-registered work include: a larger Java sample to
determine whether the RQ4/RQ5 coverage trend reaches significance; a multi-iteration
(rather than single-shot) repair loop with an explicit stopping criterion tied to suite size,
to test whether Python's regression is specific to one-shot repair or persists with iteration
budgets; replicating the study with larger models (GPT-4o, Claude-class) to test whether the
observed limits are model-size effects; extending the dataset beyond two Python libraries and
toward higher complexity bands (CC $> 10$) to locate the breaking point more precisely; and
integrating the generation--measurement loop into CI so that LLM tests are evaluated
continuously rather than in one-off experiments.
